% ─────────────────────────────────────────────────────────────────────────────
%  Capítulo 10 – Suite de pruebas
% ─────────────────────────────────────────────────────────────────────────────
\chapter{Suite de pruebas}
\label{chap:pruebas}

\section{Organización}

Las pruebas se organizan en tres grupos principales:

\begin{table}[H]
  \centering
  \caption{Grupos de pruebas}
  \label{tab:grupos}
  \begin{tabular}{llr}
    \toprule
    \textbf{Directorio}            & \textbf{Propósito}                        & \textbf{Tests} \\
    \midrule
    \texttt{tests/backend/}        & Compilación y ejecución completa           & 5 \\
    \texttt{tests/semantic/}       & Errores semánticos esperados               & 3 \\
    \texttt{tests/verification/}   & Verificación formal con Z3                 & 3 \\
    \midrule
    Total                          &                                            & 11 \\
    \bottomrule
  \end{tabular}
\end{table}

El script \texttt{tests/run\_tests.sh} ejecuta todos los grupos y
reporta el resultado de cada prueba como \texttt{[PASS]} o \texttt{[FAIL]}.

\section{Pruebas de backend}

Cada prueba compila el fichero \texttt{.java}, ejecuta el binario
resultante y compara la salida con el valor esperado.

\begin{table}[H]
  \centering
  \caption{Pruebas de backend (compilación y ejecución)}
  \label{tab:backend}
  \begin{tabular}{lp{6cm}l}
    \toprule
    \textbf{Fichero}               & \textbf{Qué se verifica}                  & \textbf{Estado} \\
    \midrule
    \texttt{t01\_print.java}       & \texttt{System.out.println} de enteros     & PASS \\
    \texttt{t02\_arithmetic.java}  & Las cuatro operaciones, módulo y precedencia & PASS \\
    \texttt{t03\_if\_else.java}    & Condicionales con ramas vacías y anidadas  & PASS \\
    \texttt{t04\_while.java}       & Bucles, suma acumulativa, retorno          & PASS \\
    \texttt{t05\_functions.java}   & Llamadas entre métodos, paso de valores    & PASS \\
    \bottomrule
  \end{tabular}
\end{table}

\section{Pruebas semánticas}

Estos ficheros son programas \emph{incorrectos}: el compilador debe
detectar el error semántico y terminar con código de salida distinto
de cero.

\begin{table}[H]
  \centering
  \caption{Pruebas de errores semánticos}
  \label{tab:semantic}
  \begin{tabular}{lp{6.5cm}l}
    \toprule
    \textbf{Fichero}               & \textbf{Error esperado}                   & \textbf{Estado} \\
    \midrule
    \texttt{sem01\_undeclared.java} & Variable usada sin declarar              & PASS \\
    \texttt{sem02\_type\_error.java}& Asignación de \texttt{double} a \texttt{int} sin conversión & PASS \\
    \texttt{sem03\_redecl.java}     & Redeclaración de variable en el mismo ámbito & PASS \\
    \bottomrule
  \end{tabular}
\end{table}

\section{Pruebas de verificación formal}

\begin{table}[H]
  \centering
  \caption{Pruebas de verificación con Z3}
  \label{tab:verif}
  \begin{tabular}{lp{5.5cm}ll}
    \toprule
    \textbf{Fichero}               & \textbf{Qué se verifica}  & \textbf{VCs} & \textbf{Estado} \\
    \midrule
    \texttt{demo\_verified.java}   & absValue, wrongMethod, sumLoop & 7 & PASS \\
    \texttt{ver01\_simple.java}    & Método sin bucles, @post básico & 2 & PASS \\
    \texttt{ver02\_invariant.java} & Invariante de bucle con variante & 5 & PASS \\
    \bottomrule
  \end{tabular}
\end{table}

\subsection{Salida de ejemplo: \texttt{demo\_verified.java}}

\begin{verbatim}
$ ./compiler --verify tests/verification/demo_verified.java
[OK]  [absValue] correccion (unsat)
[OK]  [sumLoop] iniciacion del invariante (unsat)
[OK]  [sumLoop] preservacion del invariante (unsat)
[OK]  [sumLoop] uso del invariante (unsat)
[OK]  [sumLoop] correccion (unsat)
[OK]  [sumLoop] terminacion (unsat)
[BUG] [wrongMethod] correccion (sat) <-- bug detectado
\end{verbatim}

\section{Script de ejecución}

El script \texttt{tests/run\_tests.sh} automatiza todo el proceso:

\begin{enumerate}
  \item Compila el proyecto con \texttt{make}.
  \item Ejecuta cada prueba de backend: compila el \texttt{.java} con
    \texttt{./compiler}, ejecuta el binario y compara con
    \texttt{expected/}.
  \item Ejecuta las pruebas semánticas comprobando que el compilador
    termina con código $\neq 0$.
  \item Ejecuta las pruebas de verificación comprobando que Z3 esté
    disponible en el \texttt{PATH}.
  \item Imprime un resumen: \texttt{N/M tests passed}.
\end{enumerate}

\begin{verbatim}
$ make && tests/run_tests.sh
[PASS] t01_print
[PASS] t02_arithmetic
[PASS] t03_if_else
[PASS] t04_while
[PASS] t05_functions
[PASS] sem01_undeclared
[PASS] sem02_type_error
[PASS] sem03_redecl
[PASS] ver01_simple
[PASS] ver02_invariant
[PASS] demo_verified
11/11 tests passed
\end{verbatim}
